% !PRINTED 54
\weberpage{54}
%<<CONT

manner of representation is nevertheless not wholly indifferent, and that the first often deserves preference. Especially in those cases in which the coefficients are numbers and the number-theoretic behavior of these coefficients is at issue, one must not overlook that through the binomial coefficients a numerical element foreign to the matter is introduced. That Gauss in the theory of quadratic forms (in the Disq. ar.) uses the notation with binomial coefficients when he represents the quadratic forms by $ax^2+2bxy+cy^2$, and that this notation has come into general use, has led in number theory to an unnecessary and very regrettable complication.

\begin{center}
{\large \S.\ 14.}

\vspace{0.5em}
{\large\gesperrt{Derivative of a Product.}}
\end{center}

The derived functions which we have considered here are no others than the differential quotients known from differential calculus; however, here, where only whole rational functions are involved, we have obtained the concept without using infinitesimal calculus, from the coefficients of development of the powers of $y$ in the function $f(x+y)$. If we denote the $\nu$th derivative of $f(x)$ by $D_\nu f$, then by (2), \S~13:
\begin{equation}
 f(x+y)=f(x)+yD_1f+{y^2\over1\cdot2}D_2f+{y^3\over1\cdot2\cdot3}D_3f+\cdots,
\tag{1}
\end{equation}
and from this follow at once the two fundamental rules expressed in the formulae
\begin{equation}
 D_\nu(Cf)=C D_\nu f;
\tag{2}
\end{equation}
\begin{equation}
 D_\nu(f+\varphi)=D_\nu f+D_\nu\varphi,
\tag{3}
\end{equation}
where $C$ is a constant, and $\varphi$ a second whole rational function of $x$.

A generalization of formula (2) gives the representation of the derivatives of the product $f\varphi$. Namely, if after (1) one writes, for abbreviation,
\begin{align}
 f(x+y)&=u_0+yu_1+y^2u_2+\cdots+y^n u_n,\notag\\
 \varphi(x+y)&=v_0+yv_1+y^2v_2+\cdots+y^m v_m,
\tag{4}
\end{align}
so that
\begin{equation}
 u_\nu={D_\nu f\over\Pi(\nu)},\qquad
 v_\nu={D_\nu\varphi\over\Pi(\nu)},
\tag{5}
\end{equation}
%CONT>>
